\documentclass[10pt,a4paper]{article}
\usepackage[T1]{fontenc}
\usepackage[utf8]{inputenc}
\usepackage{lmodern}
\usepackage{amsmath,amssymb}
\usepackage[margin=21mm]{geometry}
\usepackage{microtype}
\usepackage[hidelinks]{hyperref}
\usepackage{enumitem}
\setlist[itemize]{leftmargin=*,nosep,topsep=3pt}
\setlength{\parindent}{0pt}
\setlength{\parskip}{5pt}
\setlength{\emergencystretch}{3em}
\newcommand{\question}[1]{\par\medskip{\normalsize\bfseries #1}\par\smallskip}
\pagestyle{plain}

\begin{document}
{\Large\bfseries Deep brain stimulation of symptom-specific networks in Parkinson's disease\par}

\textbf{Paper:} Rajamani et al. (2024).\quad Multicentre connectomic DBS study.\par
\href{https://doi.org/10.1038/s41467-024-48731-1}{DOI: \nolinkurl{10.1038/s41467-024-48731-1}}\par
{\footnotesize Page references below use PDF page numbers, including front matter.}

\question{To which general question does this paper relate?}
What DBS can control and how to choose the performance objective. Examines whether different symptoms benefit from stimulation of different networks.

\question{How does it answer the general questions?}
\textbf{Evidence.} Stimulation locations and tract engagement were related to changes in tremor, rigidity, bradykinesia and axial symptoms. The study included 232 retrospective patients and five additional prospective patients (PDF pp.~1--10).

\textbf{Usable insight.} Symptom improvement maps to partly distinct networks. The proposed network-blending approach weights these networks according to an individual's symptom profile when choosing stimulation settings. This supports a symptom-dependent performance objective rather than assuming one physiological feature captures all treatment benefit (PDF pp.~7--10).

\textbf{Model relevance.} Tract engagement can inform hypotheses about where control enters a network and which outputs matter. It does not directly determine the notebook's temporal gains, delays, nonlinearities, or measurement matrix. Connectomic associations and the small prospective demonstration do not establish a universal stimulation-to-symptom transfer function; normative tractography is a stated limitation.

\question{What papers does it cite that could provide more information?}
\begin{itemize}
\item Horn et al. (2017), \emph{Connectivity Predicts Deep Brain Stimulation Outcome in Parkinson Disease}: connectivity-based outcome prediction (ref.~19, PDF p.~14).
\item Roediger et al. (2023), \emph{Automated deep brain stimulation programming based on electrode location}: randomized crossover programming trial (ref.~41, PDF p.~15).
\end{itemize}

\end{document}
